\section{Introduction}
\label{sec:introduction}

Edge--cloud infrastructures increasingly host services whose computation cannot be treated as a single data-center placement problem. Industrial monitoring, interactive analytics, distributed inference, and scientific services combine short response-time requirements with heterogeneous execution sites, changing network paths, and different regional electricity mixes \cite{rahmani2025optimizing,pournazari2025computation,Etheredge2026Pulse}. A scheduler operating across this continuum must decide not only \emph{where} a task runs, but also when it should start, which performance state the selected processor should use, and whether moving task state across regions is worth the transfer cost. These decisions are coupled: a faster site can shorten a critical completion path while consuming more power; a cleaner region can require a network transfer; and temporal deferral can exploit a lower-carbon interval only when the remaining deadline slack is sufficient.

The carbon dimension makes this coupling especially important. Regional grid intensity can change substantially within the lifetime of a workload, so a placement that is energy-efficient is not necessarily carbon-efficient. Carbon-aware execution studies have shown that spatial and temporal flexibility can reduce emissions, but the attainable gain depends on workload slack, resource heterogeneity, and the carbon variation actually exposed to the scheduler \cite{West2026CarbonAwareWorkflows,son2025greenscale,zhai2025forecast}. Recent work has consequently moved from cloud-only energy minimization toward carbon-aware orchestration in the cloud-to-edge continuum, including function-variant offloading, spatio-temporal federated orchestration, and energy-harvesting edge--cloud allocation \cite{calavaro2026functionvariants,alshammaa2026federated,fu2026ehcarbon}. In latency-sensitive systems, however, carbon cannot be optimized independently of deadline feasibility and the power states activated by a placement.

Three technical difficulties dominate this scheduling regime. First, sparse and heavily loaded workloads reward different resource-activation policies. A broad active pool improves headroom under contention but may waste baseline energy when only a few tasks are present. Second, a single scalar objective is brittle when the operating regime changes. Carbon, energy, and completion time can exchange rank across candidate schedules even when all schedules meet their deadlines. Third, migration is neither uniformly harmful nor uniformly beneficial: it can expose a cleaner or faster execution site, but it adds transfer delay, network energy, and transfer-related emissions. These effects become harder to reason about when node sleep/wake transitions and DVFS are available at the same time.

Existing schedulers address parts of this problem through different mechanisms. Carbon-aware methods use forecasts or regional intensity signals to shift computation \cite{zhai2024edge,zhai2025forecast,calavaro2026functionvariants}; Pareto methods preserve several trade-off solutions instead of fixing one weight vector \cite{abraham2025multi,Etheredge2026Pulse}; adaptive metaheuristics alter search behavior as the workload changes \cite{satouf2025metaheuristic,tandon2025novel}; and reinforcement-learning schedulers learn policies from observed system states \cite{shen2025reinforcement,anand2025eadrl,nasir2026relto,sravan2026drl}. These directions are complementary, yet two gaps remain relevant to practical multi-region scheduling. The first is \emph{regime robustness}: a method that is tuned for dense contention can activate too many resources on sparse traces, whereas a conservative energy policy may lengthen the critical tail under SLA pressure. The second is \emph{accounting robustness}: if unavoidable infrastructure leakage is mixed with scheduler-controlled energy, a scheduler can appear insensitive simply because a large fixed floor dominates the reported total.

We address these issues with \method, a guarded multi-profile search framework for joint placement, temporal shifting, and DVFS control. The scheduler begins from a common dynamic-greedy reference and constructs carbon-, energy-, latency-, and hybrid-oriented candidate pools. Pool size is adapted to the workload pressure estimated from arrival windows and fastest feasible execution times. Candidate schedules are then refined under an exact objective-call budget, so the search effort is directly comparable with competing metaheuristics. A protected critical-tail stage examines the latest-finishing tasks and creates narrowly targeted turbo and fastest-feasible-node repairs. Crucially, the balanced output is not selected by unconstrained scalarization. It must satisfy fixed guards on controllable carbon, controllable energy, gross energy, gross carbon, and makespan relative to the common reference, the pre-repair anchor, and a deterministic legacy activation anchor.

A second design choice is the evaluation model. We separate the energy and emissions that can change with scheduling decisions from the infrastructure floor that remains even when nodes are in deep-off leakage. Both quantities are reported. The \emph{scheduler-controllable} layer reveals the effect of placement, timing, DVFS, wake transitions, and migration, while the \emph{gross fixed-horizon} layer preserves the complete facility-level result. This distinction avoids treating background leakage as a scheduling gain and makes the observed trade-offs interpretable across workload scales.

The experimental design is deliberately split into two evidence regimes. Twelve frozen representative scenarios retain the original Google- and Alibaba-derived workloads at 30, 60, and 90 tasks. A separate SLA-constrained track contains 12 controlled stress cells created from four 150-task base scenarios and three predeclared workload seeds. Eight methods receive the same external simulator, decision space, feasibility repair, and per-scenario objective-call budget, including NSGA-II and a PPO-based scheduler. Search seeds are summarized inside each scenario before paired statistical inference, so repeated optimizer seeds are not treated as independent workloads.

The main contributions are as follows:
\begin{itemize}
    \item We formulate a joint edge--cloud scheduling model in which node placement, bounded temporal shifting, five-state DVFS, migration, sleep/wake behavior, and regional carbon intensity are evaluated within one shared decision space.
    \item We develop a workload-adaptive multi-profile search strategy that constructs energy-, carbon-, hybrid-, and latency-oriented active pools and allocates a fixed evaluation budget across deterministic structure and stochastic refinement.
    \item We introduce protected critical-tail repair and guard-based balanced selection, which target the latest-finishing tasks without allowing latency repair to ignore predeclared energy, carbon, and gross-accounting limits.
    \item We introduce transparent two-layer accounting that reports scheduler-controllable energy and emissions separately from the fixed infrastructure floor while retaining gross fixed-horizon values for system-level comparison.
    \item We provide a reproducible evaluation protocol with eight equal-budget methods, workload-level paired statistics, Pareto indicators, ablations, sensitivity analysis, exact search-call accounting, and schedule-level artifact replay.
\end{itemize}

The remainder of the paper is organized as follows. Section~\ref{sec:related} reviews carbon-aware, adaptive, learning-based, and multi-objective schedulers. Section~\ref{sec:model} defines the system and accounting model. Section~\ref{sec:method} presents \method, followed by its algorithmic properties in Section~\ref{sec:properties}. Section~\ref{sec:experimental} describes the traces, baselines, fairness controls, and statistics. Section~\ref{sec:results} reports the results, Section~\ref{sec:discussion} interprets the observed trade-offs, and Section~\ref{sec:conclusion} concludes the paper.
